\documentclass[11pt]{article}

\usepackage[margin=1in]{geometry}
\usepackage{amsmath, amssymb, amsfonts}
\usepackage{booktabs}
\usepackage{enumitem}
\usepackage{graphicx}
\usepackage{float}
\usepackage[colorlinks=true, linkcolor=blue, urlcolor=blue, citecolor=blue]{hyperref}

\graphicspath{{figures/}}

\title{Second Experiment: Linear and Neural PPO and Monte Carlo Control on a Nonlinear Problem with an Exact Optimum}
\author{Jon Legasa}
\date{\today}

\begin{document}

\maketitle

\section*{Summary}

%% RESULTS-SUMMARY

%----------------------------------------------------------------------------------------------------
\section{Setup}
\label{sec:setup}

This is the experiment of the first run (\emph{First Experiment: PPO and Monte Carlo Control
Against the Exact LQR Optimum}) repeated on the nonlinear benchmark of the notes \emph{Measuring
Reinforcement Learning Against the Exact LQR Optimum}, Section ``A nonlinear benchmark with an
exact solution''. Everything not stated here is as in the first run.

\paragraph{Problem.} Two states, one control, $Q = I$, $R = 1$, $\gamma = 0.99$:
\begin{gather*}
  x_1^+ = 0.5\, x_1 + 0.3571\, x_2, \qquad
  x_2^+ = s(x_1)\, \big(0.9193\, x_1 - 0.6436\, x_2\big) + \beta(x_1)\, u, \\
  \beta(x_1) = \frac{2 + \cos(k x_1)}{3}, \qquad
  s(x_1) = \sqrt{\frac{1 + \gamma p_2\, \beta(x_1)^2}{1 + \gamma p_2}}, \qquad p_2 = 1.5, \qquad k = 2, \\
  c(x, u) = x_1^2 + x_2^2 + u^2, \qquad x_0 = (1, -1) .
\end{gather*}
The control acts with a strength $\beta$ between $1/3$ and $1$ that depends on $x_1$. The agents
see the state and the reward $r_t = -c(x_t, u_t)$, never the equations. There is no early
termination and the dynamics are never clipped.

\paragraph{Baseline.} The cost-to-go of this nonlinear problem is exactly the Riccati solution of
its linearization at the origin, $V(x) = x^\top P x$ with $P = \mathrm{diag}(2, 1.5)$, so
\[
  J^* = x_0^\top P x_0 = 3.5 ,
\]
and the optimal policy is
\[
  u^*(x) = -\frac{\gamma p_2\, \beta(x_1)\, s(x_1)}{1 + \gamma p_2\, \beta(x_1)^2}\,
  \big(0.9193\, x_1 - 0.6436\, x_2\big),
\]
which is not linear in $x$. Three reference policies place the algorithms:
\begin{center}
  \begin{tabular}{@{}lll@{}}
    \toprule
    Policy & Gain $K$ in $u = -Kx$ & Gap \\
    \midrule
    no control, $u = 0$ & $(0, 0)$ & $0.416$ \\
    LQR gain of the linearization & $(0.5493, -0.3846)$ & $0.0308$ \\
    best linear policy & $(0.3717, -0.3286)$ & $0.00436$ \\
    optimal policy $u^*$ & not linear & $0$ \\
    \bottomrule
  \end{tabular}
\end{center}
The best linear gain is found by a numerical search over $K$ (\texttt{envs.best\_linear\_gap}). Its
gap, $0.00436$, is a floor for every policy of the form $u = -Kx$: no amount of training lets a
linear policy go below it. A table or a neural network is not bound by it. This floor is the
question of this experiment.

\paragraph{Score.} At each checkpoint the current deterministic policy of every seed is scored by its
gap $g = 1 - J^*/J \in [0, 1]$, with $g = 1$ when the cost is infinite. Unlike the first run, every
algorithm, the linear PPO included, is scored by a deterministic rollout of 5000 steps from $x_0$
($\gamma^{5000} \approx 10^{-22}$): on a nonlinear system there is no closed form for the cost of a
gain. Linear PPO is scored by its mean action $u = -Kx$, neural PPO by its mean action
$u = \mu_\theta(x)$, Monte Carlo control by its greedy table. For the table and the network, the
\emph{effective gain} is the discounted least-squares fit $u \approx -Kx$ along that rollout.

\paragraph{Instability.} For every algorithm a run is counted as unstable when its evaluation rollout
diverges or ends farther from zero than it started (over its last 100 steps). The eigenvalue test of
the first run does not apply to a nonlinear closed loop.

\paragraph{Seeds and checkpoints.} All three algorithms use the same $M = 40$ evaluation seeds as the
first run (master seed 0), and every random number of a seed comes from its own generator. Monte
Carlo control and linear PPO are checkpointed every $5 \times 10^4$ steps. Neural PPO collects
2048 steps per update, so its 21 checkpoints fall on multiples of 2048 steps, about every
$4.9 \times 10^4$ steps. Orderings involving neural PPO compare each seed's latest checkpoint at or
before a common step, without interpolation.

\paragraph{Confidence bands and figures.} The bands are 95\% Hoeffding bands, simultaneous over 21
checkpoints, of half-width $\varepsilon = 0.290$, as in the first run. Every gap axis is logarithmic;
gaps at or below $10^{-8}$ are drawn on the dashed optimum line, and the best linear policy is drawn
as a dash-dotted line at $0.00436$.

%----------------------------------------------------------------------------------------------------
\section{Algorithms}
\label{sec:algorithms}

\subsection{Monte Carlo control}

As in the first run (Sutton and Barto, Section 5.4), with the same settings. With two states the
table has $41^2 = 1681$ state cells instead of 41, each with 61 actions. A table can hold any
feedback law up to its resolution, so it is not bound by the linear floor; its own floor, set by the
bins and the action grid, has not been computed.

\subsection{PPO with a linear-Gaussian policy}

As in the first run, with the same settings: the policy is $\mathcal{N}(-Kx, \sigma^2)$ with
$K \in \mathbb{R}^{1 \times 2}$, the value function is quadratic in $x$, and the gradients are written
out by hand. It is the reference for the floor: its best possible gap is $0.00436$.

\subsection{PPO with neural networks}

The same algorithm with the policy mean and the value function given by neural networks, following
the continuous-control setup of the PPO paper (Schulman et al.\ 2017, Section 6.1 and Table 3):
\[
  \pi_\theta(u \mid x) = \mathcal{N}\big(\mu_\theta(x),\ \sigma^2\big), \qquad V_\phi(x),
\]
with $\mu_\theta$ and $V_\phi$ two separate multilayer perceptrons and $\log \sigma$ a learned number
that does not depend on the state. The mean can represent a nonlinear feedback law, so this policy is
not bound by the linear floor. Gradients are computed by PyTorch. Each seed has its own networks and
its own random numbers, so a seed's run does not depend on the batch (this is tested).

\begin{table}[H]
  \centering
  \begin{tabular}{@{}lll@{}}
    \toprule
    Setting & Value & Source \\
    \midrule
    hidden layers & 2 of 64 units, tanh & paper \\
    policy output & mean of a Gaussian; $\sigma$ learned, state-independent & paper \\
    shared parameters & none (separate policy and value networks) & paper \\
    entropy weight $c_2$ & 0 & paper \\
    steps per update $T$ & 2048 (one environment) & paper \\
    epochs & 10 & paper \\
    minibatch size & 64 (32 minibatches per epoch) & paper \\
    Adam stepsize & $3 \times 10^{-4}$, constant & paper \\
    discount $\gamma$, GAE $\lambda$ & 0.99, 0.95 & paper \\
    clip $\varepsilon$ & 0.2 & paper (best value of their search) \\
    training length & $10^6$ steps & paper \\
    \midrule
    value weight $c_1$ & 0.5 & ours (the paper: ``irrelevant'' without sharing) \\
    initial $\sigma$ & 0.5 & ours (not stated) \\
    initialization & orthogonal; policy output layer scaled by 0.01 & ours (not stated) \\
    gradient clipping & norm 0.5 per seed, policy and value separately & ours (not stated) \\
    advantage normalization & per minibatch & ours (not stated) \\
    observation and reward scaling & none & ours (not stated) \\
    \bottomrule
  \end{tabular}
  \caption{Neural PPO settings: what the paper states, and what it leaves open.}
\end{table}

The paper states that ``each of $N$ (parallel) actors collect $T$ timesteps'' but does not give $N$
for this benchmark; one environment collecting $T = 2048$ steps is our reading. The output layer of
the policy is scaled down so that the untrained policy applies controls close to zero; the system is
stable without control, so no run diverges at the start.

%% RESULTS-SECTIONS

%----------------------------------------------------------------------------------------------------
\section{Reproducing}

From the repository root:
\begin{verbatim}
uv run main.py baseline --problem nlqr --problem-set k=2
uv run main.py crude --algos monte_carlo ppo_linear ppo --problem nlqr --problem-set k=2
\end{verbatim}
%% RUN-DIRECTORY
The figures in \texttt{figures/} are copied from that run's \texttt{report/} directory.

\end{document}
